\section{Расчётная задача 6}

Направляющий вектор~$a \in \RR^3$ с координатами~$(1, 3, 1)^\top$ задаёт одномерное подпространство~$L_1 \leq \RR^3$, а линейное однородное уравнение
\[
	5x - 2y + 2z = 0
\]
--- двумерное подпространство~$L_2 \leq \RR^3$. Требуется:
\begin{enumerate}[1)]
	\item построить матрицы~$A$ проектирования на~$L_1$ параллельно~$L_2$ и~$B$ проектирования на~$L_2$ параллельно~$L_1$;
	\item доказать, что $AB = BA = 0$, $A^2 = A$, $B^2 = B$, $A + B = E$;
	\item найти ядро, образ, собственные числа и собственные векторы каждого оператора;
	\item выяснить, на какие компоненты распадается вектор~$b$ с координатами $(1, 1, 1)^\top$ при разложении~$\RR^3 = L_1 \oplus L_2$.
\end{enumerate}
\begin{solution}
	Разберём более общую ситуацию. Пусть есть разложение
	\[
		V = L_1 \oplus L_2
		,
	\]
	где $\dim{V} = n$, $\dim{L_1} = k$, $\dim{L_2} = n - k$, и требуется построить проектор, скажем, на первое подпространство:
	\[
		\pi_1 \colon L_1 \oplus L_2 \to L_1
		.
	\]
	Пусть~$\uu = (u_1, \dots, u_k) \subset L_1$ --- базис. Так как сумма~$L_1 \oplus L_2$ --- прямая, то существует набор линейных функционалов~$\xi^1$, \dots, $\xi^k \in V^*$, такой, что
	\[
		L_2 =
		\bigcap_{i = 1}^k
			\ker{\xi^i}
		.
	\]
	(Каждый из них задаёт~$(n-1)$-мерное подпространство~$\ker{\xi^i}$, и можно выбрать~$\xi^i$ так, чтобы пересечение этих подпространств было~$(n-k)$-мерным --- причём таким, каким нужно нам.) При этом выберем набор ковекторов~$(\xi^i)$ так, чтобы они образовывали базис, дуальный базису~$\uu \subset L_1$:
	\[
		\xi^i(u_j) = \delta^i_j
		.
	\]
	Тогда проектор на~$L_1$ вдоль~$L_2$ имеет вид
	\begin{equation}
		\label{eq:num_06:gen:pi_1}
		\pi_1 =
		\sum_{i = 1}^k
			u_i \otimes \xi^i
		.
	\end{equation}
	(Как и любой линейный эндоморфизм, $\pi_1 \in V^* \otimes V$.)
	
	Запишем теперь оператор~$\pi_1$, заданный формулой~\eqref{eq:num_06:gen:pi_1}, в координатах в произвольном базисе~$\ee = (e_1, \dots, e_n) \subset V$ и дуальном ему базисе~$\ee^* = (e^1, \dots, e^n)$.

	Пусть $u_i = u_i^j e_j$, где $u_i^j \in \RR$ и используется соглашение суммирования Эйнштейна. Пусть, далее, $\xi^i = \xi^i_j e^j$, где $\xi^i_j \in \RR$. Тогда
	\begin{align*}
		\pi_1 &=
		\sum_{i = 1}^k
			u_i \otimes \xi^i
		=
		\sum_{i = 1}^k
		\sum_{\mu = 1}^n
		\sum_{\nu = 1}^n
			u_i^\mu e_\mu \otimes \xi^i_\nu e^\nu
		\nlinea{=}
		\sum_{i = 1}^k
		\sum_{\mu = 1}^n
		\sum_{\nu = 1}^n
			u_i^\mu \xi^i_\nu \cdot e_\mu \otimes e^\nu
		=
		\sum_{\mu = 1}^n
		\sum_{\nu = 1}^n
		\left( 
			\sum_{i = 1}^k
				u_i^\mu \xi^i_\nu 
		\right)
		e_\mu \otimes e^\nu
		.
		\refstepcounter{equation}
		\tag{\theequation}
		\label{eq:num_06:crd:pi_1}
	\end{align*}
	Если $\ee$ записать в строку, $\ee^*$ --- в столбик:
	\begin{align*}
		\ee &=
		\Vect{
			e_1 & \dots & e_n
		}
		,
		&
		\ee^* &=
		\Vect{
			e^1 \\ \dots \\ e^n
		}
		,
	\end{align*}
	то вектор~$u_i$ и ковектор~$\xi^j$ запишутся в виде
	\begin{align*}
		u_i
		&=
		\Vect{
			e_1 & \dots & e_n
		}
		\Vect{u_i^1 \\ \dots \\ u_i^n},
		&
		\xi^j
		&=
		\Vect{\xi^j_1 & \dots & \xi^j_n}
		\Vect{
			e^1 \\ \dots \\ e^n
		}
		,
	\end{align*}
	и тогда
	\begin{align*}
		u_i \otimes \xi^j
		&=
		\left[ 
			\Vect{
				e_1 & \dots & e_n
			}
			\Vect{u_i^1 \\ \dots \\ u_i^n}
		\right]
		\otimes
		\left[ 
			\Vect{\xi^j_1 & \dots & \xi^j_n}
			\Vect{
				e^1 \\ \dots \\ e^n
			}
		\right]
		\nlinea{=}
		\Vect{
			e_1 & \dots & e_n
		}
		\left[ 
			\Vect{u_i^1 \\ \dots \\ u_i^n}
			\Vect{\xi^j_1 & \dots & \xi^j_n}
		\right]
		\Vect{
			e^1 \\ \dots \\ e^n
		}
		% \nlinea{=}
		% \Vect{
		% 	e_1 & \cdots & e_n
		% }
		% \begin{pmatrix}
		% 	u_i^1 \xi_1^j
		% 	&
		% 	\cdots
		% 	&
		% 	u_i^1 \xi_n^j
		% 	\\
		% 	\hdotsfor{3}
		% 	\\
		% 	u_i^n \xi_n^j
		% 	&
		% 	\cdots
		% 	&
		% 	u_i^n \xi_n^j
		% \end{pmatrix}
		% \Vect{
		% 	e^1 \\ \cdots \\ e^n
		% }
	\end{align*}
	Если теперь посмотреть на то, что находится внутри суммы по~$i$ в выражении~\eqref{eq:num_06:crd:pi_1}, то станет видно, что
	\newlength{\vectheight}%
	\settoheight{\vectheight}{%
		$\displaystyle
			\Vect{
				\crd{\xi^1} \\
				\cdots \\
				\crd{\xi^n}
			}
		$%
	}%
	\begin{align*}
		\pi_1 &=
		\Vect{
			e^1 \\ \dots \\ e^n
		}
		\underbrace{
			\begin{pmatrix}
				u_1^1 & u_2^1 & \dots & u_n^1
				\\
				u_1^2 & u_2^2 & \dots & u_n^2
				\\
				\hdotsfor{4}
				\\
				u_1^n & u_2^n & \dots & u_n^n
			\end{pmatrix}
		}_{U}
		\underbrace{
			\begin{pmatrix}
				\xi_1^1 & \xi_2^1 & \dots & \xi_n^1
				\\
				\xi_1^2 & \xi_2^2 & \dots & \xi_n^2
				\\
				\hdotsfor{4}
				\\
				\xi_1^n & \xi_2^n & \dots & \xi_n^n
			\end{pmatrix}
		}_{\Xi}
		\Vect{
			e_1 & \dots & e_n
		}
		\nlinea{=}
		\Vect{
			e^1 \\ \dots \\ e^n
		}
		\left[ 
			\vphantom{\rule{0pt}{\vectheight}}
		\right.
			\underbrace{
				\vphantom{
					\Vect{
						\crd{\xi^1} \\
						\cdots \\
						\crd{\xi^n}
					}
				}
				\Vect{
					| & & | \\
					\crd{u_1} &
					\cdots &
					\crd{u_n}
					\\
					| & & | \\
				}
			}_{U}
			\underbrace{
				\Vect{
					- & \crd{\xi^1} & -\\
					& \cdots & \\
					- & \crd{\xi^n} & -
				}
			}_{\Xi}
		\left.
			\vphantom{\rule{0pt}{\vectheight}}
		\right]
		\Vect{
			e_1 & \dots & e_n
		}
		.
	\end{align*}
	Произведение матриц~$P_1 := U\Xi$ и будет матрицей оператора~$\pi_1$ в базисе~$\ee$.

	В виде
	\[
		U\Xi =
		\Vect{
			\crd{u_1} &
			\dots &
			\crd{u_n}
		}
		\Vect{
			\crd{\xi^1} \\
			\dots \\
			\crd{\xi^n}
		}
	\]
	процедуру построения проектора запомнить особенно легко, особенно вспомнив формулу~$\pi = \sum_{i=1}^k u_i \otimes \xi^i$.

	Кстати, по этой же схеме строят проекторы в квантовой механике: ковекторы записываются в бра-скобках~$\bra{\cdot}$, векторы записываются в кет-скобках~$\ket{\cdot}$ (вместе --- бра и кет, \emph{bracket}, обозначение П.\,А.\,М.~Дирака). Обозначение тензорного произведения часто опускается, а применение ковектора~$\xi^i$ к вектору~$u_j$ записывается в виде~$\braket{\xi^i}{u_j}$.
	Тогда
	\[
		\pi_1 =
		\sum_{i = 1}^k
			\ket{u_i}
			\bra{\xi^j}
		,
	\]
	и результат применения~$\pi_1$ к вектору~$v \in V$ записывается в виде
	\[
		\pi_1(v) =
		\sum_{i = 1}^k
			\ket{u_i}
			\braket{\xi^j}{v}
		.
	\]
	Сразу видно, что $\braket{\xi^j}{v} \in \RR$, а будучи умноженным на вектор~$\ket{u_i}$ оно также даёт вектор. Сумма векторов --- тоже вектор, и тогда~$\pi_1(v) \in V$.

	\paragraph{Случай в задаче.}
	В нашей задаче $V \supset \ee = (e_1, e_2, e_3)$ --- стандартный ортонормированный базис; $V = L_1 \oplus L_2$, где $L_1 = \left< a\right>$, $L_2 = \ker{\xi}$ и
	\begin{align*}
		a &=
		\ee
		\Vect{1 \\ 3 \\ 1}
		,
		&
		\xi &=
		\Vect{5 & -2 & 2}
		\ee^*
		.
	\end{align*}
	Итак, $V = \left< a\right> \oplus \ker{\xi}$.
	Так как сумма прямая, то должно быть~$\xi(a) \neq 0$:
	\[
		\xi(a) =
		\Vect{5 & -2 & 2}
		\Vect{1 \\ 3 \\ 1}
		=
		5 - 6 + 2 = 1 \neq 0.
	\]
	Тогда любой вектор~$v \in V$ единственным образом раскладывается в сумму
	\begin{equation}
		\label{eq:num_06:part:v}
		v = \lambda a + w
		,
	\end{equation}
	где~$\lambda \in \RR$, $w \in \ker{\xi}$. Найдём~$\lambda$, применив к~\eqref{eq:num_06:part:v} функицонал~$\xi$:
	\[
		\xi(v) = \xi(\lambda a) + \xi(w) = \lambda \xi(a)
		,
	\]
	так как~$w \in \ker{\xi}$. Поэтому
	\[
		\lambda = \frac{\xi(v)}{\xi(a)},
	\]
	так как~$\xi(a) \neq 0$.

	Поэтому проектор~$\pi_1$ на~$L_1 = \left< a\right>$ вдоль~$L_2$ равен
	\[
		\pi_1(v) =
		\frac{1}{\xi(a)}\,\xi(v)\,a,
	\]
	или, что то же самое,
	\[
		\pi_1(\ket{v}) =
		\frac{1}{\braket{\xi}{a}}\,\braket{\xi}{v} \ket{a},
	\]
	откуда видно, что
	\[
		\pi_1 =
		\frac{1}{\braket{\xi}{a}}\,
		\ket{a} \otimes \bra{\xi}
		.
	\]
	В базисе~$\ee$ матрица~$P_1$ проектора~$\pi_1$ будет иметь вид
	\begin{multline}
		P_1 =
		\frac{1}{\braket{\xi}{a}}
		\Vect{
			a^1 \\ a^2 \\ a^3
		}
		\Vect{
			\xi_1 & \xi_2 & \xi_3
		}
		=
		\frac{1}{1}
		\Vect{
			1 \\ 3 \\ 1
		}
		\Vect{
			5 & -2 & 2
		}
		\nline{=}
		\begin{pmatrix}
			1 \cdot 5 & 1 \cdot (-2) & 1 \cdot 2
			\\
			3 \cdot 5 & 3 \cdot (-2) & 3 \cdot 2
			\\
			1 \cdot 5 & 1 \cdot (-2) & 1 \cdot 2
		\end{pmatrix}
		=
		\begin{pmatrix}
			5 & -2 & 2
			\\
			15 & -6 & 6
			\\
			5 & -2 & 2
		\end{pmatrix}
		.
		\label{eq:num_06:P1}
	\end{multline}

	\subsection[Проектор на $L_2$]{Проектор на \mathversion{bold}$L_2$}

	Чтобы найти~$\pi_2 \colon L_1 \oplus L_2 \to L_2$, можно пойти двумя путями:
	\begin{enumerate}[1)]
		\item воспользоваться свойством~$\pi_1 + \pi_2 = \id_V$ и сказать, что~$\pi_2 = \id_V - \pi_1$, и тогда в матричной форме
			\begin{equation}
				P_2 = E - P_1 =
				\begin{pmatrix}
					-4 & 2 & -2
					\\
					-15 & 7 & -6
					\\
					-5 & 2 & -1
				\end{pmatrix}
				;
				\label{eq:num_06:P2}
			\end{equation}
		\item найти базис~$\uu = (u_1, u_2) \subset L_2$ и воспользоваться общим подходом, изложенным выше.
	\end{enumerate}
	Так как в условии задачи просят \emph{проверить} равенство~$P_1 + P_2 = E$, то мы не можем использовать его и пойти по первому пути. Нужно сделать всё честно.

	\begin{remark}
		Впрочем, мы можем \emph{доказать} свойство~$\pi_1 + \pi_2 = \id_V$; тогда равенство~$P_1 + P_2 = E$ будет следовать автоматически из свойств матричной записи операторов.
	\end{remark}
	% \begin{proof}
	% 	\color{gray}
	% 	Пусть $\pi_1 \colon L_1 \oplus L_2 \to L_1$ --- проектор. Так как~$V = L_1 \oplus L_2$, а $\ker{\pi_1} = L_2$, то $V \cong \ker{\pi_1} \oplus \im{\pi_1}$. Мы применили теорему о гомоморфизме групп:
	% 	\[
	% 		\im{\pi_1} \cong V / L_2 =
	% 		\frac{L_1 \oplus L_2}{L_1}
	% 		\cong
	% 		L_1
	% 		.
	% 	\]
	% 	Итак,
	% 	\begin{align*}
	% 		\ker{\pi_1} &= L_2,
	% 		\\
	% 		\im{\pi_1} &= L_1.
	% 	\end{align*}
	% 	Для $\pi_2 \colon L_1 \oplus L_2 \to L_2$ ситуация аналогичная:
	% 	\begin{align*}
	% 		\ker{\pi_2} &= L_1,
	% 		\\
	% 		\im{\pi_2} &= L_2.
	% 	\end{align*}
	% 	Итак,
	% 	\[
	% 		V = L_1 \oplus L_2 \cong \im{\pi_1} \oplus \im{\pi_2},
	% 	\]
	% 	так что $\pi_1 + \pi_2 = \id_V$.
	% \end{proof}
	\begin{proof}
		Возможно, более наглядный способ это увидеть такой.
		Пусть $v \in V$. Так как $V = L_1 \oplus L_2$, то существует единственное разложение
		\[
			v = u_1 + u_2,
		\]
		где $u_1 \in L_1$, $u_2 \in L_2$. При этом
		\begin{align*}
			\pi_1(v) &= u_1,
			\\
			\pi_2(v) &= u_2,
		\end{align*}
		так что
		\[
			(\pi_1 + \pi_2)(v) =
			\pi_1(v) + \pi_2(v) =
			u_1 + u_2 = v =
			\id_V (v)
		\]
		для любого вектора~$v \in V$. 
	\end{proof}


	\subsection[Базис в~$\ker{\xi}$]{\mathversion{bold}Базис в~$\ker{\xi}$}

	Найдём~$\uu$. Для этого решим уравнение~$\xi(v) = 0$. Если зафиксирован базис~$\ee$ и~$v = \ee (v^1, v^2, v^3)^{\top}$, то
	\[
		\xi(v) =
		\Vect{
			5 & -2 & 2
		}
		\Vect{
			v^1 \\ v^2 \\ v^3
		}
		=
		0
		.
	\]
	Возьмём ненулевой минор этой системы, например~$\det(5) \neq 0$. Тогда $v^2 =: c^2$ и $v^3 =: c^3$ --- две константы, и
	\[
		5v^1 = 2c^2 - 2c^3,
	\]
	откуда
	\[
		\Vect{
			v^1 \\ v^2 \\ v^3
		}
		=
		\Vect{
			(2c^2 - 2c^3)/5
			\\
			c^2
			\\
			c^3
		}
		=
		c^2
		\Vect{
			2/5
			\\
			1
			\\
			0
		}
		+
		c^2
		\Vect{
			-2/5
			\\
			0
			\\
			1
		}
		=
		5c^2
		\Vect{
			2
			\\
			5
			\\
			0
		}
		+
		5
		c^2
		\Vect{
			-2
			\\
			0
			\\
			5
		}
		.
	\]
	Видно, что если $v \in \ker{\xi}$, то $v$ представляется в виде линейной комбинации векторов
	\[
		\tu_1 =
		\ee
		\Vect{2 \\ 5 \\ 0}
		\qquad
		\text{и}
		\qquad
		\tu_2 =
		\ee
		\Vect{-2 \\ 0 \\ 5}
		.
	\]
	Тогда возьмём их нормированные версии в качестве базиса. Так как~$\norm{\tu_1}^2 = \norm{\tu_2}^2 = 29$, то
	\begin{equation}
		\label{eq:num_06:U}
		\uu :=
		\ee
		\,
		\frac{1}{\sqrt{29}}
		\underbrace{
			\left( 
				\begin{array}{r|r}
					2 & -2 \\
					5 & 0 \\
					0 & 5
				\end{array}
			\right)
		}_{U}
		.
	\end{equation}

	\begingroup
	\color{gray}
	Теперь подберём такие ковекторы~$\bra{\xi^1}$, $\bra{\xi^2}$, что
	\[
		\ker{\xi^1} \cap \ker{\xi^2} = \left< a\right>
		.
	\]
	Для этого достаточно, чтобы
	$\ker{\xi^1} \supset \left< a\right>$,
	$\ker{\xi^2} \supset \left< a\right>$.
	







	Чтобы сработала формула
	\[
		\pi_2 =
		\ket{u_1} \otimes \bra{\xi^1} +
		\ket{u_2} \otimes \bra{\xi^2}
		,
	\]
	достаточно взять ортонормированный базис~$\uu$ и такие ковекторы~$\bra{\xi^i}$, что координаты вектора и ковектора совпадают:~$\xi^i_j = u_i^j$. (В общем случае это происходит с помощью изоморфизма~$V \cong V^*$, заданного матрицей Грама~$G$.)

	\begin{align*}
		\ket{u_1} \otimes \bra{\xi^1} &=
		\frac{1}{\sqrt{29}}
		\Vect{
			2 \\ 5 \\ 0
		}
		\frac{1}{\sqrt{29}}
		\Vect{
			2 & 5 & 0
		}
		=
		\frac{1}{29}
		\begin{pmatrix}
			4 & 10 & 0 \\
			10 & 25 & 0 \\
			0 & 0 & 0
		\end{pmatrix}
		;
		\\
		\ket{u_2} \otimes \bra{\xi^2} &=
		\frac{1}{\sqrt{29}}
		\Vect{
			-2 \\ 0 \\ 5
		}
		\frac{1}{\sqrt{29}}
		\Vect{
			-2 & 0 & 5
		}
		=
		\frac{1}{29}
		\begin{pmatrix}
			4 & 0 & -10 \\
			0 & 0 & 0 \\
			-10 & 0 & 25
		\end{pmatrix}
		;
		\\
		\ket{u_1} \otimes \bra{\xi^1} +
		\ket{u_2} \otimes \bra{\xi^2} &=
		\frac{1}{29}
		\begin{pmatrix}
			8 & 10 & -10 \\
			10 & 25 & 0 \\
			-10 & 0 & 25
		\end{pmatrix}
		.
	\end{align*}
\endgroup

	\subsection{Доказательство равенств}

	Равенство~$P_1 P_2 = P_2 P_1 = 0$ означает сразу две вещи:
	\begin{enumerate}[1)]
		\item что коммутатор матриц~$[P_1, P_2] = 0$ (а значит, и коммутатор операторов равен нулю) --- то есть операторы \emph{коммутируют};
		\item что $P_1 P_2 = 0$ --- а значит, что если сначала подействовать~$\pi_2$, а затем~$\pi_1$ (или наоборот в силу коммутативности), то получим~$0$.
	\end{enumerate}
	Это значит, что $V$ раскладывается в прямую сумму~$\im{\pi_1} \oplus \im{\pi_2}$.

	Равенства~$P_1^2 = P_1$ и $P_2^2 = P_2$ означают, что если подействовать оператором~$\pi_\nu$ на вектор~$v$ два раза подряд, то это то же самое, что подействовать на него единожды%
	\footnote{%
		А ещё это означает, что пространство~$V$ с фиксированным оператором~$\pi_\nu$ можно рассмотреть как факторкольцо полиномов~$\kk[t]/(t^2 - t)$. И тут нам открываются глубины...
	}%
	. Это интуитивно ясно: подействовав на~$v$ оператором~$\pi_\nu$ в первый раз, мы его спроецировали на подпространство~$\im{\pi_\nu}$; но оператор~$\pi_\nu$ не двигает~$\im{\pi_\nu}$ (мы его так построили; это \emph{не} общий факт). Поэтому повторное применение~$\pi_\nu$ оставляет образ~$\pi_\nu(v)$ вектора~$v$ на месте.

	А равенство~$P_1 + P_2 = E$ мы уже доказали в общем случае выше.

	Давай теперь посчитаем всё вышеперичисленное в матрицах. Напомню, что~$P_1$ мы вычислили в~\eqref{eq:num_06:P1}, а~$P_2$ --- в~\eqref{eq:num_06:P2}.

	\subsubsection{Коммутаторы}
		\begin{align*}
			P_1 P_2 &=
			\left(
				\begin{array}{rrr}
					5 & -2 & 2
					\\
					\hline
					15 & -6 & 6
					\\
					\hline
					5 & -2 & 2
				\end{array}
			\right)
			\left(
				\begin{array}{r|r|r}
					-4 & 2 & -2
					\\
					-15 & 7 & -6
					\\
					-5 & 2 & -1
				\end{array}
			\right)
			\nlinea{=}
			\left(
				\begin{array}{ccc}
					-20 + 30 - 10 &
					10 - 14 + 4 &
					-10 + 18 - 2
					\\
					-60 + 90 - 30 &
					30 - 42 + 12 &
					-30 + 36 - 6
					\\
					-20 + 30 - 10 &
					10 - 14 + 4 &
					-10 + 18 - 2
				\end{array}
			\right)
			=
			\begin{pmatrix}
				0 & 0 & 0 \\
				0 & 0 & 0 \\
				0 & 0 & 0
			\end{pmatrix}
			;
			\\
			P_2 P_1 &=
			\left(
				\begin{array}{rrr}
					-4 & 2 & -2
					\\
					\hline
					-15 & 7 & -6
					\\
					\hline
					-5 & 2 & -1
				\end{array}
			\right)
			\left(
				\begin{array}{r|r|r}
					5 & -2 & 2
					\\
					15 & -6 & 6
					\\
					5 & -2 & 2
				\end{array}
			\right)
			\nlinea{=}
			\left(
				\begin{array}{ccc}
					-20 + 30 - 10 &
					8 - 12 + 4 &
					-8 + 12 - 4
					\\
					-75 + 105 - 30 &
					30 - 42 + 12 &
					-30 + 42 - 12
					\\
					-25 + 30 - 5 &
					10 - 12 + 2 &
					-10 + 12 - 2
				\end{array}
			\right)
			=
			\begin{pmatrix}
				0 & 0 & 0 \\
				0 & 0 & 0 \\
				0 & 0 & 0
			\end{pmatrix}
			.
		\end{align*}

	\subsubsection{Квадраты}
		Покажем, что~$P_\nu^2 = P_\nu$.
		\begin{align*}
			P_1^2 &=
			\left(
				\begin{array}{rrr}
					5 & -2 & 2
					\\
					\hline
					15 & -6 & 6
					\\
					\hline
					5 & -2 & 2
				\end{array}
			\right)
			\left(
				\begin{array}{r|r|r}
					5 & -2 & 2
					\\
					15 & -6 & 6
					\\
					5 & -2 & 2
				\end{array}
			\right)
			\nlinea{=}
			\left(
				\begin{array}{ccc}
					25 - 30 + 10 &
					-10 + 12 - 4 &
					10 - 12 + 4
					\\
					75 - 90 + 30 &
					-30 + 36 - 12 &
					 30 - 36 + 12
					\\
					25 - 30 + 10 &
					-10 + 12 - 4 &
					10 - 12 + 4
				\end{array}
			\right)
			=
			\left(
				\begin{array}{rrr}
					5 & -2 & 2
					\\
					15 & -6 & 6
					\\
					5 & -2 & 2
				\end{array}
			\right)
			;
			\\
			P_2^2 &=
			\left(
				\begin{array}{rrr}
					-4 & 2 & -2
					\\
					\hline
					-15 & 7 & -6
					\\
					\hline
					-5 & 2 & -1
				\end{array}
			\right)
			\left(
				\begin{array}{r|r|r}
					-4 & 2 & -2
					\\
					-15 & 7 & -6
					\\
					-5 & 2 & -1
				\end{array}
			\right)
			\nlinea{=}
			\left(
				\begin{array}{ccc}
					16 - 30 + 10 &
					-8 + 14 - 4 &
					8 -12 + 2
					\\
					60 - 105 + 30 &
					-30 + 49 - 12 &
					30 - 42 + 6
					\\
					20 - 30 + 5 &
					-10 + 14 - 2 &
					10 - 12 + 1
				\end{array}
			\right)
			=
			\left(
				\begin{array}{rrr}
					-4 & 2 & -2
					\\
					-15 & 7 & -6
					\\
					-5 & 2 & -1
				\end{array}
			\right)
			.
		\end{align*}

	\subsection{Ядра, образы и спектры}
		Как мы уже замечали выше,
		\begin{align*}
			\ker{\pi_1} &= L_2,
			&
			\ker{\pi_2} &= L_1,
			\\
			\im{\pi_1} &= L_1,
			&
			\im{\pi_2} &= L_2,
		\end{align*}
		причём
		\begin{align*}
			L_1 &= \left< a \right>,
			\\
			L_2 &= \Ann{\xi} = \ee U,
		\end{align*}
		где~$U$ --- матрица из~\eqref{eq:num_06:U}.

		При этом базис в~$L_1$ --- это в точности единственный собственный вектор оператора~$\pi_1$, а базис~$\uu$ --- это базис в собственном подпространстве~$L_2$, так как
		\begin{align*}
			\pi_1 L_1 &\subset L_1,
			\\
			\pi_2 L_2 &\subset L_2,
		\end{align*}
		а это и характеризует собстенные подпространства.

		Впрочем, можно посчитать это и руками.

		% \subsubsection[Оператор $\pi_1$]{\mathversion{bold}Оператор $\pi_1$}

	\subsection{Образ вектора}
		Образы вектора
		\[
			b := \ee \Vect{1 \\ 1 \\ 1} = \ee \crd{b}
		\]
		при разложении~$V = L_1 \oplus L_2$ --- это в точности $\pi_1(b)$ и~$\pi_2(b)$:
		\begin{align*}
			\pi_1(b) &= P_1 \crd{b}
			=
			\left(
				\begin{array}{rrr}
					5 & -2 & 2
					\\
					15 & -6 & 6
					\\
					5 & -2 & 2
				\end{array}
			\right)
			\Vect{1 \\ 1 \\ 1}
			=
			\left(
				\begin{array}{r}
					5
					\\
					15
					\\
					5
				\end{array}
			\right)
			,
			\\
			\pi_2(b) &= P_2 \crd{b}
			=
			\left(
				\begin{array}{rrr}
					-4 & 2 & -2
					\\
					-15 & 7 & -6
					\\
					-5 & 2 & -1
				\end{array}
			\right)
			\Vect{1 \\ 1 \\ 1}
			=
			\left(
				\begin{array}{r}
					-4
					\\
					-14
					\\
					-4
				\end{array}
			\right)
			.
		\end{align*}
\end{solution}
